\section{Conclusion}

The empirical exact-rung result does reveal a theorem-level design principle,
but the principle is more specific than ``finish with $\omega=1$.''  Optimal
SOR creates a known reflection debt.  Acceleration requires delaying that
debt while useful mass propagates; locality requires preventing the delayed
state and neighbor-generated backflow from repeatedly scanning the processed
region.  The correct ladder spacing is $1+\Theta(\sqrt\alpha)$, which yields
the desired $1/\sqrt\alpha$ number of epochs.

The new block theorem shows that exact debt resolution is already local on
forests and width-sensitive on cyclic regions.  In particular, the spider's
triangular inward wave is not a lower bound for elimination-aware methods.
It also shows that an exact block rung deletes all relaxation history on a
fixed region, so the only possible accelerated effect of the SOR ladder is
faster free-boundary discovery.
The regularization-path ledger then proves that output-sensitive discovery
would give $O((w+1)^2/(\rho\sqrt\alpha))$ total work with no additional
dynamic-range logarithm.  The support-radius estimate is graph-universal, and
fixing the seed value gives a monotone scalar homotopy on every graph, so
cycles do not create a correctness or support-discovery obstruction.  On
trees the online work theorem is complete:
piecewise-affine child responses expose exact activation events, lazy
iteration never enters an inactive descendant, and Chebyshev inverse decay
bounds the event-propagation depth.  The result is an exact, condition-free
$\widetilde O(1/(\rho\sqrt\alpha))$ local solver on every finite tree.  Paths
and symmetric spiders admit the even stronger output-linear final-$\rho$
solution and therefore beat the persistent-support product charge.
Bounded-size biconnected cyclic gadgets also meet the product scale through
lazy block responses.  Their block identifiers and final block--cut tree do
not need to be supplied: zero-valued activation events make online mergers
prefix-preserving, so affected response iterators restart only from the
current checkpoint.  Large blocks are no longer uniformly open: radial
Schur elimination solves a whole cycle in output-linear work.  More strongly,
orthogonal shell averaging gives the same output-linear exact solver on every
root-distance-equitable graph, including thick-shell families such as
hypercubes and distance-regular graphs.  The symmetry test is certifying on an
arbitrary graph and falls back safely at the first irregular event.  Still
more generally, online refinement plus lazy responses on a hidden equitable
tree quotient permit several inequivalent thick cells per distance layer and
achieve the
exact $\widetilde O(1/(\rho\sqrt\alpha))$ product bound even on original
graphs with unbounded treewidth and large biconnected cores, without supplied
cell identifiers.  Independently,
the exact boundary gate reaches the product scale whenever only
$\chi_\star$ support vertices
occupy each distance shell; this includes fixed-width cyclic strips.  The
approximate-output target is weaker still.  Certified one-sided intervals
suffice for support-safe admission and termination, scalar response leverage
is uniformly at most one, and exterior demand changes are energy contractions.
Source-aware sensitivity has an exact scalar state: it is the loss of terminal
Schur diagonal under the current batch elimination, and these losses telescope
per exterior label.  A margin-adaptive hierarchy may use additive error at its
current pruning scale.  Chebyshev inverse-square-root probes build a complete
one-sided boundary-loss snapshot in
$\widetilde O(\vol(S)/\sqrt\alpha)$ work, so all full rebuilds at geometric
anchors already fit the product ledger.  Source normalization inside an epoch
also needs no exact frontier square root: a constant-factor spectral frontier
and Chebyshev recurrence prepare every Gaussian source probe in
$\widetilde O(T_{\rm mv}/\sqrt\alpha)$ work.  The positive resolvent theorem
goes further: the inverse square root is a constant-factor positive sum of
shifted exact rungs whose total work mass and monotone Schur-coordinate
settlement count have the $1/\sqrt\alpha$ dependence.  Each shifted rung is a
sparse face system, and every unfinished local solve is represented exactly
by nonnegative anchor debt; exact cleanup produces precisely the shifted Schur
residual.  Thus only energy-packed interval crossings created by light updates
between checkpoints need to be reported.
Small cut rank closes the target-side harmonic bank, while a comb construction
refutes any universal low-rank reduction.  The remaining obstruction is now
online graph-work scheduling on high-rank cores: the companion response note
already flushes every fixed-face aggregate debt at the product scale, but one
must locate all intervening crossings by partial output-sensitive work without
forcing that full flush or a boundary scan after every light admission inside
a large core with neither bounded
articulation blocks, an equitable tree quotient, nor thin radial support.
Exact trace compression or
an exact-gate batch bound remain stronger alternatives.  The gate does
dominate proximal gradient and therefore has an exact margin-sensitive
termination bound, but a three-vertex path proves that its per-batch objective
contraction can be only $1-\Theta(\alpha)$.  Consequently the missing
$1/\sqrt\alpha$ argument, if true, must be combinatorial or data-structural:
it cannot be recovered by a black-box strongly-convex energy estimate.

\bibliographystyle{plainnat}
\bibliography{../../references}
